\section{A Data-Driven Approach}\label{sec:data-driven}






In this section, we present the main contribution of this work which provides a data-driven approach to calibrating \POFUL{}. Note that $\iota_t$ and $\tau_t$ are parameters of \POFUL{} that can be controlled by a decision-maker, the essential point is to find a computable, non-trivial upper bound $\hat{\alpha}_t$ for the uncertainty ratio $\alpha_t$, which turns into an upper bound $\hat{\mu}_t$ for the regret proxy $\mu_t$ that's deeply related to the frequentist regret of \POFUL{}. 




% \blue{Typically, we can check that $\hat{\alpha}_t$ is no less than one because when $ \tilde{x}_t = x_t^\star$, we have $\alpha_t = \left\|x_t^{\star}\right\|_{V_t^{-1}} /\left\|\widetilde{x}_t\right\|_{V_t^{-1}} = 1$ and $\hat{\alpha}_t$ is an upper bound for $\alpha_t$. As a result, it holds that $\hat{\mu}_t = (1+\hat{\alpha}_t)(1+\iota_t) +  (1-\hat{\alpha}_t) \tau_t \leq  (1+\hat{\alpha}_t)(1+\iota_t) $, i.e., setting $\tau_t  = 0$ yields a valid upper bound for   $\hat{\mu}_t$. Given this observation, \blue{we will} focus on scenarios where $\tau_t = 0$ for all $t\in [T]$.}  

We focus on scenarios where $\tau_t = 0$ for all $t\in [T]$. These include LinTS and variants such as TS-Freq, as well as Greedy - standard algorithms still lacking theoretical regret guarantees. Below, we construct upper bounds $\{\hat{\alpha}_t\}_{t\in[T]}$ for the continuous-action scenario. Bounds for discrete-action scenarios appear in Appendix \ref{sec:discrete_actions}.
 
 
 
 















\subsection{Continuous action space}
\label{sec:bound_sst_ratio}  

%When it comes to the generic settings where the number of actions is large or when the set of actions is a continuous set in $\RR^d$, the upper bound above is inefficient or even infeasible to compute. In this section, we tackle this issue. 

Our strategy capitalizes on geometric insights related to the properties of the confidence ellipsoids, providing upper bounds that can be computed efficiently.  For the sake of a better illustration, we consider 
 $\cX_t = \cS_{d-1}$ for all $t\in[T]$ for this scenario, where $\cS_{d-1} =\left\{x \in \mathbb{R}^d:\|x\|=1\right\}$ is the unit hypersphere in $\mathbb{R}^d$. This is a standard example of continuous action space, and is the same as the setting considered in \cite{abeille2017linear}. We remark that for this specific setting, the problem is still hard. This is because we don't have a closed-form solution for the set of potentially optimal actions.
 
 In this setting, the optimal action $x_t^\star(\theta )  \coloneqq \argmax_{x\in \cX_t} \dotp{x}{\theta}$ takes the form $ x^\star_t(\theta ) = \theta/\|\theta\|$. To upper bound $\alpha_t$, we consider respectively the smallest and largest value of $\|x^\star_t(\theta)\|_{V_t^{-1}}$ for $\theta$ in the confidence ellipsoids of $\theta$, namely, $\cE^{\RLS}_{t,\delta^\prime, \lambdareg}$ and $\cE^{\PVT}_{t,\delta^\prime, \lambdareg}$. Specifically, we have 
  \# \label{eq:alpha_hat_raw}
 \alpha_t \leq \frac{\sup _{\theta \in \cE_{t,\delta^\prime, \lambdareg}^{\RLS}}\|x^\star_t(\theta)\|_{V_t^{-1}}}{ \inf _{\theta \in \cE_{t,\delta^\prime, \lambdareg}^{\PVT}}\|x^\star_t(\theta)\|_{V_t^{-1}}}.
 \#
 

 
As is illustrated in Figure~\ref{fig:ellipsoid1}, the set of potentially optimal actions $\cC_t$ is the projection of the confidence ellipsoid $\cE_t$  onto $\cS_{d-1}$. It's hard to get a closed-form expression for $\cC_t$, so we cannot directly calculate the range of $V_t^{-1}$-norm of actions in $\cC_t$. Nevertheless, when \POFUL{} has implemented sufficient exploration so that $\cE_t$ is small enough, $\cC_t$ concentrates accordingly to a small cap on $\cS_{d-1}$. Therefore, it is possible to estimate the range of the $V_t$-norm by employing geometric reasoning. Subsequently, this estimated range will be utilized to ascertain the range of the $V_t^{-1}$-norm.

% All actions within $\cC_t$  point to similar directions and thus have similar $V_t$-norm. (Note that for a unit vector $x$, we have that $\|x\|_{V_t}$ can be written as a weighted summation of the eigenvalues of $V_t$, where the weights are determined by the direction of $x$.) 




% Hence we can bound the range of $V_t$-norm by investigating the concentration of  $\cC_t$, and we'll use it to determine the range of $V_t^{-1}$-norm.
 
 
  

 

 \begin{figure}[htp]
	\centering
	\begin{subfigure}[b]{0.45\textwidth}
	\includegraphics[width=0.9\textwidth]{figures/ellipsoid_8} 
		\caption{}
	\end{subfigure}
 %\hfill
\begin{subfigure}[b]{0.45\textwidth}
	\includegraphics[width=0.9\textwidth]{figures/ellipsoid_4} 
		\caption{}
	\end{subfigure}
	\caption{Illustration of potentially optimal actions set $\cC_t$ in $\RR^2$. (a): $\cC_t$ is $\cE_t$'s projection onto $\cS_{d-1}$. (b): As more data is collected, $\cE_t$ shrinks (colors show exploration levels). Potentially optimal actions point in similar directions, determining their $V_t$-norm. This suggests their $V_t$-norm range could be estimated geometrically. }
	\label{fig:ellipsoid1}
\end{figure}
%
The main theorem (proved in Section \ref{sec:proof_sketch_alpha}) derives an upper bound for $\alpha_t$ based on this idea.




\begin{theorem}\label{thm:alpha}


Suppose $\cX_t = \cS_{d-1}$ for all $t \in [T]$. Define $m_t = (\|\hat{\theta}_t\|_{V_t}^2  - (\beta^{RLS}_{t,\delta^\prime, \lambdareg})^2)/(\|\hat{\theta}_t\| +  \beta^{RLS}_{t,\delta^\prime, \lambdareg}/\lambda_d(V_t) )^2$, $ M_t = \|\hparam_t\|_{V_t^2}^2/(\|\hat{\theta}_t\|_{V_t}^2  - (\beta^{\PVT}_{t,\delta^\prime, \lambdareg})^2)$. Let  $k\in[d]$ be the integer that satisfies $\lambda_{k}(V) \leq M_t \leq \lambda_{k+1}(V) $.
 Define $\beta^{\PVT}_{t,\delta^\prime, \lambdareg} \coloneqq \iota_t \beta^{RLS}_{t,\delta^\prime, \lambdareg}$ and
					\$
	\Phi_t &=   \begin{cases}
				( \lambda_{1}^{-1}(V_t) + \lambda_{d}^{-1}(V_t)   - m_t\lambda_{1}^{-1}(V_t)\lambda_{d}^{-1} (V_t))^{\frac{1}{2}}, &\ \ \ \ \textit{if}\  \|\hat{\theta}_t\|_{V_t} \geq \beta^{RLS}_{t,\delta^\prime, \lambdareg}	\\
				\lambda_d^{-\frac{1}{2}}(V_t), & \ \ \ \ \textit{if}\  \|\hat{\theta}_t\|_{V_t} <  \beta^{RLS}_{t,\delta^\prime, \lambdareg} \end{cases},\\
	\Psi_t &=   \begin{cases}
				(\lambda_{k}^{-1}(V_t) +  \lambda_{k+1}^{-1}(V_t)   - M_t\lambda_{k}^{-1}(V_t)\lambda_{k+1}^{-1} (V_t))^{\frac{1}{2}}, &\textit{if}\  \|\hat{\theta}_t\|_{V_t} \geq \beta^{\PVT}_{t,\delta^\prime, \lambdareg}	\\
				\lambda_1^{-\frac{1}{2}}(V_t), & \textit{if}\  \|\hat{\theta}_t\|_{V_t} <  \beta^{\PVT}_{t,\delta^\prime, \lambdareg}\end{cases}.
				\$			

	Then for all $t \in [T]$, conditioned on $\hat{\cA}_t \cap \tilde{\cA}_t$, it holds for all $s \leq t$ that $\alpha_s\leq \halpha_s\coloneqq \Phi_s/\Psi_s$.
 
 
 % \$
	% \inf _{\theta \in \cE_s^{\PVT}}\|x^\star_s(\theta)\|_{V_s^{-1}} \geq \Psi_s,  \quad   \sup _{\theta \in \cE_s^{RLS}}\|x^\star_s(\theta)\|_{V_s^{-1}} \leq \Phi_s , \quad \alpha_s\leq \halpha_s\coloneqq \Phi_s/\Psi_s.\$

 
\end{theorem}
To better understand what Theorem~\ref{thm:alpha} implies, we discuss some special cases in Appendix~\ref{sec:Cases} and provide empirical validations for them. 



% \begin{remark}

% Clearly, one can further improve the bound in Equation \eqref{eq:alpha_hat_raw} by plugging in the selected action $\tilde{x}_t$. This is true for standard algorithms such as OFUL, LinTS, and Greedy. Explicitly, the uncertainty ratio can be bounded using $\tilde{x}_t$ and Theorem~\ref{thm:alpha} as 
% \#
% \alpha_t \leq \frac{\sup _{\theta \in \cE_t^{\RLS}}\|x^\star_t(\theta)\|_{V_t^{-1}}}{\|\tilde{x}_t\|_{V_t^{-1}}}\leq \frac{\Phi_t}{\|\tilde{x}_t\|_{V_t^{-1}}}.
% \#


% However, for the variants of \POFUL{} that we are going to propose in the Section~\ref{sec:meta}, it is necessary to calculate the bound \emph{before} selecting the action $\tilde{x}_t$. Therefore, the bound shouldn't depend on $\tilde{x}_t$. 



% \end{remark}

\section{A Meta-Algorithm for Course-Correction}
\label{sec:meta}


This section demonstrates how the data-driven regret bound can enhance standard bandit algorithms. We propose a meta-algorithm that creates course-corrected variants of base algorithms, achieving minimax-optimal frequentist regret guarantees while preserving most original characteristics, including computational efficiency and typically low regret.


We take LinTS as an example of the base algorithm, and propose the algorithm Linear Thompson Sampling with Maximum Regret (Proxy) (TS-MR). The idea is to measure the performance of LinTS using $\hat{\mu}_t$ and avoid bad LinTS actions by switching to OFUL actions. Specifically, at each time $t$, TS-MR calculates the upper bound $\hat{\mu}_t$ and compares it with a preset threshold $\mu$. If  $\hat{\mu}_t > \mu$, LinTS might be problematic and  TS-MR takes an OFUL action to ensure a low instantaneous regret; if $\hat{\mu}_t\leq \mu$, TS-MR  takes the LinTS action. We remark that setting $\iota_t = 0$ for all $t\in [T]$ yields the corresponding Greedy-MR algorithm. The pseudocode is presented in Algorithm~\ref{alg:TS-MR} in Appendix~\ref{sec:TS-MR}. 

\begin{remark}
Computing $\hat{\mu}_t$ primarily requires SVD decomposition of the sample covariance matrix $V_t$. Since $V_t=\lambda_{\text{reg}} I_d+\sum_{s=1}^t x_s x_s^{\top}$ is updated via rank-one matrices, its SVD can be efficiently updated \citep{gandhi2017updating}, preventing computational bottlenecks.
\end{remark}


By design, course-corrected algorithms ensure that $\mu_t \leq \max\{\mu,2\}$ for all $t \in [T]$. Substituting this upper bound into Theorem~\ref{thm:oracle_regret} establishes that these algorithms achieve optimal frequentist regret, up to a constant factor.
\begin{corollary} TS-MR and Greedy-MR achieve a frequentist regret of
$\tilde{\cO}( \max\{\mu,2\} d \sqrt{T})$.
\end{corollary}

% \blue{
% \proof{} Note that $\mu_t \leq \max\{\mu,2\}$, by Theorem ~\ref{thm:oracle_regret}, we have 
% \#
% \cR(T) \leq   \sqrt{2 d  T(\max\{\mu,2\})^2\log \left(1+\frac{T}{\lambdareg}\right) 
%  }\beta_{T,\delta^\prime,\lambdareg}^{RLS} = \tilde{\cO}( \max\{\mu,2\} d \sqrt{T}).\nonumber
% \#
%  \endproof
%  }

In high-risk settings where LinTS and Greedy may fail, a small $\mu$ ensures TS-MR and Greedy-MR select more OFUL actions, promoting sufficient exploration. Conversely, in low-risk settings where the original algorithms perform well, a large $\mu$ favors TS and greedy actions, minimizing unnecessary exploration and reducing computational cost. In Appendix~\ref{sec:tune_mu}, we show how $\mu$ impacts the fraction of OFUL actions in TS-MR and Greedy-MR and their performance. Results indicate that course-corrected algorithms maintain robustness across a range of moderate $\mu$ values, suggesting that the precise selection of $\mu$ is unlikely to present a significant practical concern.

% In high-risk settings where LinTS and Greedy may fail, a small $\mu$ ensures TS-MR and Greedy-MR select more OFUL actions, promoting sufficient exploration. Conversely, in low-risk settings where the original algorithms perform well, a large $\mu$ favors TS and greedy actions, minimizing unnecessary exploration and reducing computational cost. In Appendix~\ref{sec:tune_mu}, we illustrate the impact of $\mu$ on the fraction of OFUL actions invoked in TS-MR and Greedy-MR and the performance of these two algorithms. The simulation results demonstrate that the performance of the course-corrected algorithms remains relatively robust to the choice of $\mu$. Nevertheless, a recommended heuristic for selecting $\mu$ is to ensure adequate exploration through OFUL actions during the initial stages.


% In practice, the choice of the threshold $\mu$ depends on the problem settings. In high-risk settings where LinTS and Greedy fail with high probability, one can set a small $\mu$ so that  TS-MR and Greedy-MR select more OFUL actions, guaranteeing the necessary amount of exploration. In low-risk settings where original LinTS and Greedy algorithms work well,  a large $\mu$ allows TS-MR and Greedy-MR to favor TS and greedy actions, avoiding unnecessary exploration and reducing computational cost. 







